\section{রিপোর্ট: তথ্যের উপস্থাপন}
\begin{definitionbox}{রিপোর্ট (Report)}
টেবিল বা কুয়েরির তথ্যকে নির্দিষ্ট বিন্যাসে দেখানো, সংক্ষিপ্ত করা ও মুদ্রণের উপযোগী উপস্থাপনকে \textbf{রিপোর্ট (Report)} বলে।
\end{definitionbox}
পরীক্ষার ফলপত্র, মাসিক বিক্রয়, বেতন বিবরণী ও ঠিকানার লেবেল রিপোর্টের উদাহরণ। ফর্মের প্রধান ভূমিকা এন্ট্রি ও পারস্পরিক ব্যবহার; রিপোর্টের প্রধান ভূমিকা তথ্য প্রকাশ।

\subsection{রিপোর্টের উৎস ও প্রকারভেদ}
রিপোর্টের Record Source একটি টেবিল বা কুয়েরি হতে পারে। Student ও Result-এর তথ্য একসঙ্গে দরকার হলে আগে সঠিক সংযোগের কুয়েরি তৈরি করে তার ফল রিপোর্টে দেওয়া যায়। রিপোর্টে নাম, নম্বর, মোট, গড়, শ্রেণিভিত্তিক দল ও পৃষ্ঠা নম্বর রাখা যায়। এতে ব্যবহারকারী মূল রেকর্ডের দীর্ঘ তালিকা না দেখেও প্রয়োজনীয় চিত্র বুঝতে পারে।

\textbf{Detail Report}-এ পৃথক রেকর্ডের বিস্তারিত তথ্য থাকে; যেমন প্রত্যেক শিক্ষার্থীর নম্বর। \textbf{Summary Report}-এ পৃথক সব রেকর্ডের বদলে মোট, গড় বা দলভিত্তিক সারাংশ থাকে; যেমন প্রতি বিভাগের গড় GPA। \textbf{Cross-tabulation Report}-এ দুটি বৈশিষ্ট্যের সারি-কলামভিত্তিক সারাংশ দেখানো হয়; যেমন বিভাগ ও পরীক্ষাভিত্তিক উত্তীর্ণের সংখ্যা।

\textbf{Graphics/Charts-সহ Report} চার্ট বা চিত্র দিয়ে তুলনা সহজ করে; যেমন মাসভিত্তিক বিক্রয়ের বারচিত্র। \textbf{Form/Letter Report} নির্দিষ্ট চিঠি বা ফরমের কাঠামোয় ব্যক্তিভেদে মান বসায়, যেমন শিক্ষার্থীভিত্তিক নোটিশ। \textbf{Label Report} একই মাপের লেবেলে নাম-ঠিকানা মুদ্রণ করতে ব্যবহৃত হয়। Formatted Report বলতে নির্ধারিত ফন্ট, কলাম, শিরোনাম ও অন্যান্য বিন্যাসযুক্ত প্রতিবেদন বোঝানো হয়।

\subsection{রিপোর্টের ব্যান্ড বা সেকশন}
রিপোর্টের অংশগুলো কখন ও কতবার দেখা যায় তা জানা জরুরি। পুরো রিপোর্টের শিরোনাম ও প্রতিটি পৃষ্ঠার শিরোনামের কাজ এক নয়। নিচের পার্থক্যগুলো থেকে উপযুক্ত সেকশন নির্বাচন করা যায়।
\begin{flushleft}\small
\begin{tabularx}{\linewidth}{|p{34mm}|X|}\hline
\chSixHeader ব্যান্ড & অবস্থান ও উদাহরণ\\\hline
Report Header & রিপোর্টের শুরুতে; প্রতিষ্ঠানের নাম ও প্রতিবেদনের শিরোনাম\\\hline
Page Header & প্রতিটি পৃষ্ঠার উপরে; কলামের শিরোনাম\\\hline
Detail & প্রতিটি রেকর্ডের জন্য; নাম ও নম্বরের সারি\\\hline
Page Footer & প্রতিটি পৃষ্ঠার নিচে; পৃষ্ঠা নম্বর\\\hline
Report Footer & রিপোর্টের শেষে; সর্বমোট বা সামগ্রিক সারাংশ\\\hline
Group Header/Footer & দল থাকলে দলের শুরু/শেষে; বিভাগ ও বিভাগের মোট\\\hline
\end{tabularx}
\end{flushleft}
উদাহরণস্বরূপ, ১৫ পৃষ্ঠার ফলাফলে কলেজের বড় শিরোনাম প্রথমে একবার, Roll--Name--Marks প্রতিটি পৃষ্ঠায় এবং শেষ পৃষ্ঠায় সামগ্রিক গড় রাখা যায়। মোট নম্বর প্রতিটি Detail-এ না দিয়ে Report Footer-এ রাখলে তা সমগ্র রিপোর্টের সারাংশ বোঝায়। দলভিত্তিক গড় চাইলে যথাযথ Group Footer ব্যবহার করা যায়।

\subsection{Report Wizard ও বিন্যাস}
Access-এর Create অংশে Report Wizard খুলে উৎস এবং প্রয়োজনীয় ফিল্ড নির্বাচন করি। প্রয়োজন হলে Grouping, Sorting এবং layout নির্ধারণ করি। নাম দিয়ে Finish করার পর Preview-তে পৃষ্ঠা, কলামের প্রস্থ, শিরোনাম ও মোটের স্থান যাচাই করি। যেমন StudentResult কুয়েরি থেকে Roll, Name ও ICTMark নিয়ে বিভাগভিত্তিক ফল তৈরি করা যায়।

রিপোর্টে তথ্য বেশি হলে উপযুক্ত orientation, কাগজের মাপ ও margin নির্বাচন করতে হয়। অকারণে বড় ফন্ট বা অতিরিক্ত ফাঁক দিলে পৃষ্ঠা বাড়ে; আবার অতিরিক্ত ছোট করলে পাঠযোগ্যতা নষ্ট হয়। সঠিক column width, group এবং page break ব্যবহার করে নাম বা নম্বর কেটে যাচ্ছে কি না যাচাই করতে হয়। Report View-তে দেখা এবং Print Preview-তে মুদ্রিত পৃষ্ঠার বিন্যাস পরীক্ষা পৃথক উদ্দেশ্যে কাজে লাগে।

\begin{flushleft}\small
\begin{tabularx}{\linewidth}{|p{26mm}|X|X|}\hline
\chSixHeader বিষয় & ফর্ম & রিপোর্ট\\\hline
প্রধান উদ্দেশ্য & এন্ট্রি, দেখা ও সম্পাদনা & বিন্যস্ত প্রকাশ ও মুদ্রণ\\\hline
ব্যবহারকারীর কাজ & কন্ট্রোলে মান দেওয়া/বাছাই & ফল ও সারাংশ পড়া\\\hline
উদাহরণ & ভর্তি ফর্ম & ফলাফল বিবরণী\\\hline
উৎস & টেবিল বা কুয়েরি & টেবিল বা কুয়েরি\\\hline
\end{tabularx}
\end{flushleft}
\subsection{Detail থেকে Summary: একটি পূর্ণ উদাহরণ}
ধরা যাক Result কুয়েরিতে Science বিভাগের Rupa-এর 90, Tania-এর 80 এবং Humanities বিভাগের Nabil-এর 70 নম্বর আছে। Detail Report-এ তিন শিক্ষার্থীর তিনটি সারি থাকবে। Summary Report-এ Science-এর গড় 85 ও Humanities-এর গড় 70 দেখানো যেতে পারে। সারাংশ তৈরি মানে মূল রেকর্ড মুছে ফেলা নয়; একই উৎসের ভিন্ন উপস্থাপন।
\begin{flushleft}\small
\begin{tabular}{|l|l|l|}\hline
\chSixHeader বিভাগ & নাম & নম্বর\\\hline
Science & Rupa & 90\\\hline
Science & Tania & 80\\\hline
Humanities & Nabil & 70\\\hline
\end{tabular}\qquad
\begin{tabular}{|l|l|l|}\hline
\chSixHeader বিভাগ & সংখ্যা & গড়\\\hline
Science & 2 & 85\\\hline
Humanities & 1 & 70\\\hline
\end{tabular}
\end{flushleft}
প্রত্যেক শিক্ষার্থীর নম্বর দরকার হলে শুধু সারাংশ যথেষ্ট নয়। আবার প্রতিষ্ঠানের সামগ্রিক অবস্থা জানতে দীর্ঘ তালিকার চেয়ে সারাংশ উপযোগী হতে পারে। তাই রিপোর্টের ধরন নির্বাচন করতে পাঠকের প্রশ্ন ও প্রয়োজন বুঝতে হয়।

\subsection{গ্রুপিং, সর্টিং ও পৃষ্ঠার অংশ}
Group by বিভাগ দিলে একই বিভাগের রেকর্ড একসঙ্গে থাকে। বিভাগের ভেতরে নম্বর বড় থেকে ছোট সাজানো আর বিভাগ তৈরি করা আলাদা কাজ। Group Footer-এ বিভাগের মোট/গড় এবং Report Footer-এ সমগ্র প্রতিবেদনের মোট/গড় রাখা যায়। দুই বিভাগের গড় সরাসরি যোগ করে দুই দিয়ে ভাগ করলে সবসময় সব শিক্ষার্থীর প্রকৃত গড় পাওয়া যায় না, কারণ বিভাগের শিক্ষার্থীসংখ্যা আলাদা হতে পারে।

উপরের উদাহরণে সব নম্বরের গড় $(90+80+70)/3=80$। দুই বিভাগের গড়ের সরল গড় $(85+70)/2=77.5$, যা এখানে প্রতিষ্ঠানের শিক্ষার্থীভিত্তিক গড় নয়। রিপোর্টে গণনার একক ও grouping ভুল হলে দেখতে সঠিক বিন্যাসেও ভুল তথ্য প্রকাশ হতে পারে।

\textbf{রিপোর্ট পরীক্ষা:} Preview-তে শিরোনাম ও পৃষ্ঠা নম্বর আছে কি না, দীর্ঘ নাম কেটে যাচ্ছে কি না, সব প্রয়োজনীয় রেকর্ড এসেছে কি না এবং মোটের মান উৎসের সঙ্গে মিলে কি না দেখো। Detail Band হলো রিপোর্টের একটি অংশ; Detail Report হলো রিপোর্টের একটি ধরন। শব্দ দুটি কাছাকাছি হলেও একই বিষয় নয়।

\chSixInterface{ch06_access_report.png}{Report ও Print Preview}{Detail অংশে রেকর্ড এবং Header ও Footer অংশে শিরোনাম ও সারসংক্ষেপের মতো তথ্য থাকে।}
